\documentclass{article}
\usepackage[acronym]{glossaries}

\makeglossaries

% Define acronyms
\newacronym{gcd}{GCD}{Greatest Common Divisor}
\newacronym{lcm}{LCM}{Least Common Multiple}
\newacronym{api}{API}{Application Programming Interface}
\newacronym{html}{HTML}{Hypertext Markup Language}
\newacronym{css}{CSS}{Cascading Style Sheets}

% Define glossary entries (non-acronym terms)
\newglossaryentry{latex}{
	name={LaTeX},
	description={A document preparation system used for high-quality typesetting}
}

\newglossaryentry{compiler}{
	name={compiler},
	description={A program that translates source code into executable code}
}

\newglossaryentry{algorithm}{
	name={algorithm},
	description={A step-by-step procedure for solving a problem or performing a task}
}

\newglossaryentry{variable}{
	name={variable},
	description={A named storage location in memory that holds a value}
}

\begin{document}
	
	\section*{Introduction}
	
	Given a set of numbers, there are elementary methods to compute 
	its \acrlong{gcd}, which is abbreviated \acrshort{gcd}. This process 
	is similar to that used for the \acrfull{lcm}.
	
	In computer science, an \acrfull{api} allows different software 
	applications to communicate with each other. Web development 
	often involves \acrshort{html} and \acrshort{css} for building 
	user interfaces.
	
	The \gls{latex} system uses a \gls{compiler} to process documents. 
	An \gls{algorithm} typically uses \glspl{variable} to store 
	intermediate results.
	
	\clearpage
	
	\section*{Acronyms}
	\printglossary[type=\acronymtype, title=List of Acronyms]
	
	\clearpage
	
	\section*{Glossary}
	\printglossary[title=General Glossary]
	
\end{document}